% =====================================================================
%  IRTC 工作量说明（中文 · 5 页）
%  与中文手册共用同一套视觉设计：scripts/manual-pdf/style.tex
% =====================================================================
\documentclass[11pt, oneside]{book}

\usepackage{xcolor}
\usepackage[a4paper, top=3cm, bottom=2.5cm, left=2.4cm, right=2.4cm,
            headsep=0.55cm, footskip=1.1cm]{geometry}
\usepackage{fontspec}
\usepackage{xeCJK}
\usepackage{fancyvrb}
\usepackage{longtable}
\usepackage{booktabs}
\usepackage{array}
\usepackage{hyperref}

\setmainfont{Palatino}
\setmonofont{Menlo}[Scale=0.86, HyphenChar=None]
\setCJKmainfont{Songti SC}[BoldFont=Songti SC Bold, ItalicFont=Kaiti SC]

% =====================================================================
%  IRTC 中文使用手册 —— PDF 视觉设计（XeLaTeX）
%  惟安数据科技（北京）有限公司
% =====================================================================

% ---------------------------------------------------------------- 调色板
\definecolor{IRTCink}{HTML}{102A43}      % 主色：深靛蓝
\definecolor{IRTCinkdeep}{HTML}{081726}  % 封面底色
\definecolor{IRTCinkmid}{HTML}{16375A}   % 封面渐变
\definecolor{IRTCgold}{HTML}{C8912F}     % 强调：铜金
\definecolor{IRTCgoldsoft}{HTML}{E3B65C}
\definecolor{IRTCteal}{HTML}{1B7A8C}     % 次强调：青
\definecolor{IRTCslate}{HTML}{51637A}    % 辅助文字
\definecolor{IRTCmist}{HTML}{EDF2F7}     % 浅底
\definecolor{IRTCpaper}{HTML}{F7F4EE}
\definecolor{IRTCrule}{HTML}{D5DEE7}
\definecolor{IRTCcodebg}{HTML}{F5F7FA}
\definecolor{IRTCcodeframe}{HTML}{DEE5EC}
\definecolor{IRTCcodeink}{HTML}{16506B}

% ---------------------------------------------------------------- 宏包
\usepackage{tikz}
\usetikzlibrary{calc}
\usepackage{eso-pic}
\usepackage{tcolorbox}
\tcbuselibrary{skins,breakable}
\usepackage[explicit]{titlesec}
\usepackage{titletoc}
\usepackage{fancyhdr}
\usepackage{enumitem}
\usepackage{colortbl}
\usepackage{newunicodechar}
\usepackage{ragged2e}
\usepackage{needspace}

% ---------------------------------------------------------------- 字体
% 正文：宋体 + Palatino；标题：苹方 + Avenir Next；代码：Menlo
\newfontfamily\LatinDisplay{Avenir Next}[
  UprightFont = * Medium, BoldFont = * Demi Bold, ItalicFont = * Italic]
\newCJKfontfamily\CJKdisplay{PingFang SC}[BoldFont = PingFang SC Semibold]
\newCJKfontfamily\CJKlight{PingFang SC}[UprightFont = PingFang SC Light,
  BoldFont = PingFang SC Regular]
\newcommand{\HeadFont}{\LatinDisplay\CJKdisplay}
\newcommand{\LightFont}{\LatinDisplay\CJKlight}

% 缺字回退：箭头与带圈数字取自宋体
\newfontfamily\SymFallback{Songti SC}
\newunicodechar{→}{{\SymFallback\char"2192}}
\newunicodechar{↔}{{\SymFallback\char"2194}}
\newunicodechar{①}{{\SymFallback\char"2460}}
\newunicodechar{②}{{\SymFallback\char"2461}}
\newunicodechar{③}{{\SymFallback\char"2462}}
\newunicodechar{④}{{\SymFallback\char"2463}}
\newunicodechar{⑤}{{\SymFallback\char"2464}}
\newunicodechar{⑥}{{\SymFallback\char"2465}}
\newunicodechar{⑦}{{\SymFallback\char"2466}}
\newunicodechar{⑧}{{\SymFallback\char"2467}}
\newunicodechar{⑨}{{\SymFallback\char"2468}}

% ---------------------------------------------------------------- 版面
\linespread{1.22}
\setlength{\parindent}{0pt}
\setlength{\parskip}{0.5\baselineskip}
\setlist[itemize]{leftmargin=1.35em, itemsep=2pt, topsep=3pt, parsep=0pt}
\setlist[enumerate]{leftmargin=1.6em, itemsep=2pt, topsep=3pt, parsep=0pt}
\emergencystretch=1.5em
\sloppy
\widowpenalty=10000
\clubpenalty=10000

% 行内代码着色
\let\IRTColdtexttt\texttt
\renewcommand{\texttt}[1]{{\color{IRTCcodeink}\IRTColdtexttt{#1}}}

% ---------------------------------------------------------------- 章首页
% 章首整页顶部为深色横幅：右侧压一个巨大的浅色章号，横幅内是章标签与章标题。
% 标签与标题都画在 TikZ 覆盖层里，位置与横幅严格对齐，正文另起于横幅之下。
\newcommand{\IRTCchapterhead}[3]{%
  \AddToShipoutPictureBG*{%
    \begin{tikzpicture}[remember picture, overlay]
      \fill[IRTCink] (current page.north west)
        rectangle ([yshift=-6.2cm]current page.north east);
      \fill[IRTCgold] ([yshift=-6.2cm]current page.north west)
        rectangle ([yshift=-6.32cm]current page.north east);
      \node[anchor=north east, text=white!17!IRTCink,
            font=\LatinDisplay\fontsize{92}{92}\selectfont]
        at ([xshift=-1.5cm, yshift=-0.55cm]current page.north east) {#1};
      \node[anchor=north west, text=IRTCgoldsoft,
            font=\HeadFont\fontsize{11.5}{14}\selectfont]
        at ([xshift=2.4cm, yshift=-2.95cm]current page.north west) {#2};
      \node[anchor=north west, text=white, text width=13.4cm, align=flush left,
            font=\HeadFont\fontsize{23}{31}\selectfont]
        at ([xshift=2.38cm, yshift=-3.75cm]current page.north west) {#3};
    \end{tikzpicture}%
  }%
}

\titleformat{\chapter}[display]{}{}{0pt}
  {\IRTCchapterhead{\thechapter}{第\,\thechapter\,章}{#1}%
   \markboth{第\,\thechapter\,章\quad #1}{}}
\titleformat{name=\chapter, numberless}[display]{}{}{0pt}
  {\IRTCchapterhead{}{IRTC 使用手册 · 中文版}{#1}\markboth{#1}{}}
\titlespacing*{\chapter}{0pt}{0pt}{86pt}
\titlespacing*{name=\chapter, numberless}{0pt}{0pt}{86pt}

% ---------------------------------------------------------------- 各级标题
\titleformat{\section}
  {\HeadFont\fontsize{15}{20}\selectfont\color{IRTCink}\raggedright}
  {}{0pt}
  {\textcolor{IRTCgold}{\rule[-0.05em]{3pt}{0.92em}}\hspace{0.6em}#1}
\titlespacing*{\section}{0pt}{20pt plus 4pt}{7pt}

\titleformat{\subsection}
  {\HeadFont\fontsize{12.5}{17}\selectfont\color{IRTCteal}\raggedright}
  {}{0pt}{#1}
\titlespacing*{\subsection}{0pt}{14pt plus 3pt}{5pt}

\titleformat{\subsubsection}
  {\HeadFont\fontsize{11.5}{15}\selectfont\color{IRTCslate}\raggedright}
  {}{0pt}{#1}
\titlespacing*{\subsubsection}{0pt}{11pt plus 2pt}{4pt}

% ---------------------------------------------------------------- 页眉页脚
\newcommand{\IRTCpagepill}{%
  \tikz[baseline=(x.base)]{%
    \node[fill=IRTCink, text=white, rounded corners=2.5pt,
          inner xsep=8pt, inner ysep=3.5pt,
          font=\LatinDisplay\fontsize{9}{9}\selectfont] (x) {\thepage};}}

\pagestyle{fancy}
\fancyhf{}
\renewcommand{\headrulewidth}{0.4pt}
\renewcommand{\headrule}{\ifdim\headrulewidth>0pt
  {\color{IRTCrule}\hrule height \headrulewidth}\fi}
\fancyhead[L]{\LightFont\fontsize{8.5}{10}\selectfont\color{IRTCslate}\leftmark}
\fancyhead[R]{\LightFont\fontsize{8.5}{10}\selectfont\color{IRTCslate}%
  IRTC 使用手册 · V1.1.1}
\fancyfoot[C]{\IRTCpagepill}

\fancypagestyle{plain}{%
  \fancyhf{}%
  \renewcommand{\headrulewidth}{0pt}%
  \fancyfoot[C]{\IRTCpagepill}}

\fancypagestyle{IRTCblank}{%
  \fancyhf{}\renewcommand{\headrulewidth}{0pt}}

% ---------------------------------------------------------------- 目录
\renewcommand{\contentsname}{目录}
\setcounter{tocdepth}{1}
\titlecontents{chapter}[0em]
  {\addvspace{11pt}\HeadFont\fontsize{12}{15}\selectfont\color{IRTCink}}
  {\makebox[2.1em][l]{\color{IRTCgold}\thecontentslabel}}
  {}
  {\hfill\normalfont\LatinDisplay\fontsize{10}{12}\selectfont
   \color{IRTCslate}\thecontentspage}
\titlecontents{section}[2.1em]
  {\addvspace{1.5pt}\LightFont\fontsize{9.5}{12}\selectfont\color{IRTCslate}}
  {}{}
  {\titlerule*[0.6pc]{\color{IRTCrule}.}%
   \normalfont\LatinDisplay\fontsize{9}{11}\selectfont\color{IRTCslate}\thecontentspage}

% ---------------------------------------------------------------- 引用框
\renewenvironment{quote}
  {\begin{tcolorbox}[breakable, enhanced, colback=IRTCmist, colframe=IRTCmist,
     boxrule=0pt, arc=1.5pt, left=11pt, right=11pt, top=8pt, bottom=8pt,
     borderline west={2.8pt}{0pt}{IRTCteal},
     before skip=9pt, after skip=11pt]\small}
  {\end{tcolorbox}}

% ---------------------------------------------------------------- 代码块
\makeatletter
\AtBeginDocument{%
  \@ifundefined{Shaded}{}{%
    \renewenvironment{Shaded}
      {\begin{tcolorbox}[breakable, enhanced, colback=IRTCcodebg,
         colframe=IRTCcodeframe, boxrule=0.5pt, arc=2.5pt,
         left=9pt, right=9pt, top=7pt, bottom=7pt,
         borderline west={2.2pt}{0pt}{IRTCgold},
         before skip=9pt, after skip=11pt]\small}
      {\end{tcolorbox}}}%
}
\makeatother

\DefineVerbatimEnvironment{verbatim}{Verbatim}{%
  fontsize=\small, frame=leftline, framerule=2.2pt,
  rulecolor=\color{IRTCteal}, xleftmargin=11pt, framesep=9pt,
  formatcom=\color{IRTCink}}

% ---------------------------------------------------------------- 表格
\renewcommand{\arraystretch}{1.32}
\arrayrulecolor{IRTCrule}
\setlength{\arrayrulewidth}{0.5pt}

% ---------------------------------------------------------------- 尾页
\newcommand{\IRTCendpage}{%
  \cleardoublepage
  \thispagestyle{IRTCblank}%
  \AddToShipoutPictureBG*{%
    \begin{tikzpicture}[remember picture, overlay]
      \shade[top color=IRTCinkdeep, bottom color=IRTCinkmid]
        (current page.south west) rectangle (current page.north east);
      \foreach \i in {0,1,2}{
        \draw[IRTCgoldsoft!18!IRTCinkmid, line width=0.9pt]
          ([xshift=2cm, yshift=4.2cm+\i*0.9cm]current page.south west)
          .. controls ([xshift=8cm, yshift=4.2cm+\i*0.9cm]current page.south west)
             and ([xshift=10cm, yshift=7.4cm+\i*0.9cm]current page.south west)
          .. ([xshift=18cm, yshift=7.6cm+\i*0.9cm]current page.south west);}
      \fill[IRTCgold] ([xshift=2.4cm, yshift=-6.4cm]current page.north west)
        rectangle ([xshift=4.0cm, yshift=-6.47cm]current page.north west);
    \end{tikzpicture}}%
  \vspace*{5.4cm}
  {\LatinDisplay\fontsize{40}{44}\selectfont\color{white}IRTC}\par
  \vspace{0.55cm}
  {\HeadFont\fontsize{15}{20}\selectfont\color{white}%
   项目反应理论分析工具包}\par
  \vspace{0.25cm}
  {\LightFont\fontsize{10.5}{15}\selectfont\color{IRTCgoldsoft}%
   Marginal Maximum Likelihood Estimation for Item Response Models}\par
  \vspace{1.5cm}
  {\LightFont\fontsize{10.5}{18}\selectfont\color{white}%
   版本 1.1.1（2026-07-17）· 已发布于 CRAN\par
   制作人 / 维护人：马崑翔（Kunxiang Ma）\par
   makunxiang@weiandata.com\par
   版权所有 \copyright{} 2026 惟安数据科技（北京）有限公司\par
   许可证 GPL (>= 2)\par}
  \vspace{1.2cm}
  {\LightFont\fontsize{10}{16}\selectfont\color{IRTCgoldsoft}%
   项目主页 \quad https://github.com/weiandata/IRTC\par
   CRAN \quad\ \  https://CRAN.R-project.org/package=IRTC\par
   问题反馈 \quad https://github.com/weiandata/IRTC/issues\par}
  \vfill
  {\LightFont\fontsize{9.5}{14}\selectfont\color{white}%
   本手册随软件版本更新；如手册与包内帮助页（\IRTColdtexttt{?irtc}）不一致，以帮助页为准。\par}
  \vspace{1.2cm}
  \clearpage}

\usetikzlibrary{arrows.meta}

\hypersetup{
  pdftitle={IRTC 工作量说明（版本 1.1.1）},
  pdfauthor={马崑翔 Kunxiang Ma},
  pdfsubject={IRTC 项目反应理论分析包 —— 代码量、文字量、技术路径与实现效果},
  pdfcreator={XeLaTeX},
  colorlinks=true, linkcolor=IRTCteal, urlcolor=IRTCteal}

% 页眉页脚：本文件不使用手册的页眉文案
\fancyhead[R]{\LightFont\fontsize{8.5}{10}\selectfont\color{IRTCslate}%
  IRTC 工作量说明 · V1.1.1}

% 分节页：与手册章首页同一形制
\newcommand{\WLsection}[3]{%
  \clearpage
  \thispagestyle{plain}%
  \IRTCchapterhead{#1}{#2}{#3}%
  \markboth{#3}{}%
  \vspace*{78pt}%
}

\newcommand{\WLnote}[1]{%
  \begin{tcolorbox}[enhanced, colback=IRTCmist, colframe=IRTCmist, boxrule=0pt,
    arc=2pt, left=12pt, right=12pt, top=9pt, bottom=9pt,
    borderline west={2.8pt}{0pt}{IRTCteal}, before skip=10pt, after skip=6pt]
  {\fontsize{10}{16.5}\selectfont\color{IRTCink}#1}
  \end{tcolorbox}}

\begin{document}

% =====================================================================
%  第 1 页 · 封面
% =====================================================================
\thispagestyle{IRTCblank}
\AddToShipoutPictureBG*{%
  \begin{tikzpicture}[remember picture, overlay]
    \shade[top color=IRTCinkdeep, bottom color=IRTCinkmid]
      (current page.south west) rectangle (current page.north east);
    \begin{scope}
      \clip (current page.south west) rectangle (current page.north east);
      \begin{scope}[shift={([xshift=14.6cm, yshift=0.75cm]current page.south west)},
                    x=1.35cm, y=3.3cm]
        \foreach \slope/\tint in {1.1/IRTCgoldsoft!58!IRTCinkmid,
                                  1.9/IRTCteal!62!IRTCinkmid,
                                  0.7/white!30!IRTCinkmid,
                                  1.5/IRTCgoldsoft!26!IRTCinkmid}{
          \draw[\tint, line width=1.1pt]
            plot[domain=-2.4:3.6, samples=100, smooth]
            (\x, {1/(1+exp(-\slope*(\x+0.35)))});}
      \end{scope}
    \end{scope}
    \fill[IRTCgold] ([xshift=2.6cm, yshift=-2.5cm]current page.north west)
      rectangle ([xshift=4.4cm, yshift=-2.58cm]current page.north west);
    \fill[IRTCgold!88!IRTCinkdeep] (current page.south west)
      rectangle ([yshift=0.32cm]current page.south east);
  \end{tikzpicture}}

\vspace*{1.05cm}
{\LightFont\fontsize{10.5}{14}\selectfont\color{IRTCgoldsoft}%
 R 语言项目反应理论（IRT）分析包}\par

\vspace{1.05cm}
{\LatinDisplay\fontsize{54}{56}\selectfont\color{white}IRTC}\par
\vspace{0.3cm}
{\HeadFont\fontsize{27}{34}\selectfont\color{white}工作量说明}\par
\vspace{0.42cm}
{\color{IRTCgold}\rule{3.2cm}{1.6pt}}\par
\vspace{0.42cm}
{\LightFont\fontsize{11.5}{18}\selectfont\color{IRTCgoldsoft}%
 代码量 · 文字量 · 技术路径 · 实现效果 \quad|\quad 版本 1.1.1，已发布于 CRAN}\par

\vspace{0.75cm}
{\HeadFont\fontsize{12.5}{17}\selectfont\color{white}%
 它和现在常见的 IRT 软件包有什么不一样？}\par
\vspace{0.35cm}
{\LightFont\fontsize{10}{16}\selectfont\color{white}
常用的 IRT 包（如 mirt、TAM、ltm）默认使用者是会写程序的统计人员：数据得先自己
整理成一张干净的数字表，跑完还得自己把结果从模型里抠出来、再手工排版成报告。
IRTC 把这一头一尾也接上了：\par
\vspace{0.34cm}
\hspace{0.2cm}\textcolor{IRTCgoldsoft}{\textbf{1}}\quad
\textbf{原始表格直接进}——Excel、CSV、SPSS、Stata、SAS 都能直接读，
学号列、缺失码、A/B/C/D 选项作答由软件自动识别与计分。\par
\vspace{0.16cm}
\hspace{0.2cm}\textcolor{IRTCgoldsoft}{\textbf{2}}\quad
\textbf{成品结果直接出}——只要一条命令，就生成三个可交付的 Excel 表和三种版式的
Word/HTML 报告，不用再自己排版。\par
\vspace{0.16cm}
\hspace{0.2cm}\textcolor{IRTCgoldsoft}{\textbf{3}}\quad
\textbf{结论用大白话讲}——每道题好不好、该不该改，都给中文判断、原因和建议，
而不是只给一堆系数。\par
\vspace{0.16cm}
\hspace{0.2cm}\textcolor{IRTCgoldsoft}{\textbf{4}}\quad
\textbf{算得更大}——自研的有界内存算法，让普通笔记本也能处理百万人规模的多维模型。\par
\vspace{0.16cm}
\hspace{0.2cm}\textcolor{IRTCgoldsoft}{\textbf{5}}\quad
\textbf{给自动化和 AI 留了门}——结果是结构固定的机器可读数据，报错自带错误码与修复建议。\par}

\vfill

{\color{white!35}\rule{5.5cm}{0.5pt}}\par
\vspace{0.38cm}
{\LightFont\fontsize{12}{19}\selectfont\color{white}%
 \textbf{R 包创建人 / 维护人：马崑翔}\par}
{\LightFont\fontsize{10.5}{17}\selectfont\color{IRTCgoldsoft}%
 Kunxiang Ma\par}
\vspace{0.4cm}
{\LightFont\fontsize{10.5}{16}\selectfont\color{white}%
 软件版本 1.1.1（2026-07-17）\quad · \quad 本说明 2026 年 8 月\par}
\vspace{0.7cm}

% =====================================================================
%  第 2 页 · 代码量
% =====================================================================
\WLsection{01}{第一部分}{代码量}

以仓库当前（版本 1.1.1）内容统计，项目共 \textbf{247 个源文件、19{,}874 行}。
其中真正的实现代码（R 与 C++）约 1.18 万行，自动化测试 5{,}282 行，
测试与实现的行数比约为 \textbf{1 : 2.2}。

\vspace{2pt}
{\fontsize{10}{15}\selectfont
\begin{tabular}{@{}>{\raggedright\arraybackslash}p{4.3cm}rr>{\raggedright\arraybackslash}p{6.6cm}@{}}
\toprule
\strut\textcolor{IRTCink}{\bfseries 组成部分} &
\textcolor{IRTCink}{\bfseries 文件} &
\textcolor{IRTCink}{\bfseries 行数} &
\textcolor{IRTCink}{\bfseries 内容} \\
\midrule
\strut R 源码 \IRTColdtexttt{R/} & 131 & 10{,}579 &
估计层、数据层、统计层、输出层、AI 接口与结构化报错 \\
\strut C++ 源码 \IRTColdtexttt{src/} & 11 & 1{,}255 &
并行 E 步、概率与充分统计量、2PL/GPCM 求解、流式引擎 \\
\strut 自动化测试 \IRTColdtexttt{tests/} & 64 & 5{,}282 &
309 个测试块、1{,}087 处断言 \\
\strut 帮助文档 \IRTColdtexttt{man/} & 21 & 1{,}171 &
每个导出函数的 R 标准帮助页与示例 \\
\strut 机器接口 \IRTColdtexttt{inst/} & 4 & 316 &
\IRTColdtexttt{llms.txt} API 摘要、版权声明等 \\
\strut 构建与排版 \IRTColdtexttt{scripts/} & 15 & 1{,}204 &
发布校验、基准测试、手册 PDF 设计与生成 \\
\strut 示例 \IRTColdtexttt{examples/} & 1 & 67 &
端到端使用示例 \\
\midrule
\strut\textbf{合计} & \textbf{247} & \textbf{19{,}874} & \\
\bottomrule
\end{tabular}\par}

\vspace{10pt}
\textbf{去掉空行与纯注释后的有效代码行}：R 层 8{,}479 行、C++ 层 952 行、
测试 4{,}420 行、帮助文档 1{,}063 行。

\vspace{4pt}
\textbf{对外接口规模}：20 个导出函数 + 12 个 S3 方法（\IRTColdtexttt{summary}、
\IRTColdtexttt{print}、\IRTColdtexttt{plot}、\IRTColdtexttt{anova}、
\IRTColdtexttt{logLik} 等），17 个注册的 C++ 原生符号，3 个内置仿真数据集
（\IRTColdtexttt{data.sim.rasch}、\IRTColdtexttt{data.gpcm}、\IRTColdtexttt{data.mc}）。

\WLnote{\textbf{这些行数意味着什么？}\quad
估计核心（\IRTColdtexttt{irtc.mml} 系列 + C++）是「把模型算对、算快」的部分；
而占比更大的 R 层，是把「原始表格 → 清洗 → 计分 → 体检 → 估计 → 评级 →
成品交付」这条链路补齐的部分。后者正是普通使用者能少写代码、
少犯错的原因，也是本项目区别于现有 IRT 包的主要投入所在。}

% =====================================================================
%  第 3 页 · 文字量
% =====================================================================
\WLsection{02}{第二部分}{文字量}

项目的文字工作分三类：面向使用者的手册、R 内置帮助页、以及仓库文档
（发布记录、贡献与安全说明、合规与提交记录等）。
全部 Markdown 文档合计 \textbf{21 个文件、7{,}822 行、约 29.1 万字符}，
R 帮助页另有 1{,}171 行。

\vspace{2pt}
{\fontsize{10}{15}\selectfont
\begin{tabular}{@{}>{\raggedright\arraybackslash}p{4.9cm}>{\raggedright\arraybackslash}p{4.4cm}>{\raggedright\arraybackslash}p{6.0cm}@{}}
\toprule
\strut\textcolor{IRTCink}{\bfseries 文档} &
\textcolor{IRTCink}{\bfseries 规模} &
\textcolor{IRTCink}{\bfseries 说明} \\
\midrule
\strut 中文使用手册 & 1{,}767 行 \par 63{,}262 字符（含 22{,}771 个汉字）&
15 章 + 4 个附录，从「电脑上什么都没装」讲到读懂每一个数字 \\
\strut 英文使用手册 & 675 行 \par 36{,}770 字符 &
面向英文使用者的完整手册 \\
\strut R 内置帮助页 & 21 个 Rd 文件 \par 1{,}171 行 &
\IRTColdtexttt{?irtc}、\IRTColdtexttt{?irtc.mml} 等，随包安装 \\
\strut 仓库文档 & 21 个 Markdown 文件 \par 7{,}822 行 / 29.1 万字符 &
中英文 README、版本记录、发布与提交记录、贡献与安全说明 \\
\strut AI 接口说明 & 98 行 &
\IRTColdtexttt{inst/llms.txt}，供 AI 助手一次读入的 API 与数据结构摘要 \\
\bottomrule
\end{tabular}\par}

\vspace{10pt}
除静态文档外，软件\textbf{运行时输出的文字同样是双语的}：白话摘要、
四级题目质量评级的原因与建议、数据体检的问题清单、三种受众的报告正文、
以及全部报错信息，都各有中英两套文案，随
\IRTColdtexttt{options(irtc.lang)} 切换。这部分文案写在 R 源码里，
已计入前一页的代码行数。

\vspace{4pt}
本次 1.1.1 正式发布还新增了中文手册的 PDF 排版工程（封面、引言、目录、
章节页、尾页与配色体系），共 5 个设计文件，
已计入 \IRTColdtexttt{scripts/} 的行数中。

\WLnote{\textbf{为什么写这么多字？}\quad
本项目的目标读者里有相当一部分\textbf{没有编程和统计基础}。
对这类使用者来说，「软件能算」只解决了一半问题，
「看得懂、照着做、出了错知道怎么修」才是另一半。
因此手册用零基础语言重写了全部概念，报错也全部配了原因与修复建议。}

% =====================================================================
%  第 4 页 · 技术路径
% =====================================================================
\WLsection{03}{第三部分}{技术路径}

整体结构是「两端加一芯」：数据层把各种原始表格变成可估计的数据，
估计核心用边际极大似然（MML）+ EM 算法求解，输出层把结果变成人和机器都能直接用的东西。

\vspace{6pt}
\begin{center}
\begin{tikzpicture}[font=\LightFont\fontsize{9.5}{13}\selectfont,
  box/.style={rounded corners=3pt, draw=IRTCrule, fill=IRTCmist,
              text width=3.8cm, align=center, inner ysep=9pt, text=IRTCink},
  core/.style={rounded corners=3pt, draw=IRTCink, fill=IRTCink,
               text width=5.1cm, align=center, inner ysep=10pt, text=white},
  arrow/.style={-{Straight Barb[length=4pt]}, draw=IRTCgold, line width=1pt}]
  \node[box] (a) at (0,0) {\textbf{数据层}\\[2pt] 读取 · 清洗 · 计分\\ Q 矩阵 · 数据体检};
  \node[core] (b) at (6.0,0) {\textbf{估计核心（MML + EM）}\\[2pt]
        C++ 并行 E 步 · 维度分解\\ grid / streaming 自动路由\\ SQUAREM 加速 · 可选受控精度};
  \node[box] (c) at (12.0,0) {\textbf{输出层}\\[2pt] 白话摘要 · 质量评级\\ Excel · 报告 · 图 · JSON};
  \draw[arrow] (a) -- (b);
  \draw[arrow] (b) -- (c);
  \node[text=IRTCslate, font=\LightFont\fontsize{8.5}{11}\selectfont]
    at (6.0,-1.8) {全程双语输出，全部报错带错误码 · 原因 · 修复建议};
\end{tikzpicture}
\end{center}

\vspace{4pt}
{\fontsize{10}{15}\selectfont
\begin{tabular}{@{}>{\raggedright\arraybackslash}p{4.9cm}>{\raggedright\arraybackslash}p{10.4cm}@{}}
\toprule
\strut\textcolor{IRTCink}{\bfseries 关键技术} &
\textcolor{IRTCink}{\bfseries 作用} \\
\midrule
\strut C++ 多线程与循环重排 &
把最耗时的 E 步交给 C++ 并行计算，减少重复读取与中间复制 \\
\strut 维度分解（题间多维）&
利用「每题只属于一个维度」，把主要计算量从约 N×I×Q\textsuperscript{D}
降为 N×I×Q + N×D×Q\textsuperscript{D} \\
\strut 分块流式 E 步 &
不保存全部被试的完整后验矩阵，内存由随人数增长改为受块大小控制 \\
\strut SQUAREM 与安全化 Newton &
加快 EM 收敛，并在坏初值或病态数据下保持稳定 \\
\strut 引擎自动路由 &
按人数、题数、维度与内存预测在 grid / streaming 间自动选择，并说明理由 \\
\strut 受控精度剪枝（可选）&
删除低权重积分节点换取提速，同时用分层抽样\textbf{实测}误差并出具报告 \\
\bottomrule
\end{tabular}\par}

\vspace{8pt}
\textbf{精确优先原则}：默认路径不做任何近似——完整积分网格、精确 E 步；
提速手段一律是可选开关，且开启后必须以实测误差报告为准。
阻尼、线搜索与正定修复只用于稳定求解，不改变统计目标；
任何会改变估计结果的设置都需要使用者显式开启。

% =====================================================================
%  第 5 页 · 实现效果
% =====================================================================
\WLsection{04}{第四部分}{实现效果}

{\fontsize{10}{15}\selectfont
\begin{tabular}{@{}>{\raggedright\arraybackslash}p{3.6cm}>{\raggedright\arraybackslash}p{11.7cm}@{}}
\toprule
\strut\textcolor{IRTCink}{\bfseries 方面} &
\textcolor{IRTCink}{\bfseries 已达到的结果} \\
\midrule
\strut 发布状态 &
版本 1.1.1 已通过 CRAN 审核并上架，使用者一条
\IRTColdtexttt{install.packages("IRTC")} 即可安装；Windows 与 macOS 直接获得
编译好的二进制包 \\
\strut 质量验证 &
\IRTColdtexttt{R CMD check -{}-as-cran} 0 错误 0 警告；1{,}191 项自动化测试全部通过；
整体代码覆盖率 96.0\%，关键模块 ≥ 95\%；macOS / Windows / Linux 三平台持续集成 \\
\strut 规模与性能 &
百万样本、多维 GPCM 原型测试中，流式引擎约 99 秒/次迭代，峰值内存约 2.4 GB；
同规模的传统完整后验方案可能需要 TB 级内存 \\
\strut 模型覆盖 &
Rasch/1PL、PCM、PCM2、RSM、2PL、GPCM；单维与题间多维；
潜在回归、多组比较、抽样权重、固定参数 \\
\strut 交付效果 &
一条命令从 Excel 到成品：三个可直接交付的 Excel（题目质量表、跨年度链接参数表、
被试能力表）+ 三种受众版式的 Word/HTML 报告 + 诊断图 \\
\strut 自动化友好 &
\IRTColdtexttt{irtc\_results()} / \IRTColdtexttt{irtc\_json()} 输出结构冻结的
机器可读结果；全部报错带错误码、原因与修复建议，可被程序直接捕获处理 \\
\bottomrule
\end{tabular}\par}

\vspace{12pt}
{\HeadFont\fontsize{13}{18}\selectfont\color{IRTCink}%
\textcolor{IRTCgold}{\rule[-0.05em]{3pt}{0.92em}}\hspace{0.6em}一句话总结}\par
\vspace{4pt}
同类软件解决的是「模型能不能算出来」；本项目在此之上，
把「原始数据怎么进来」「结果怎么读懂」「成品怎么出去」「机器怎么调用」
四件事一并做完，并用 CRAN 的公开审核、96\% 的测试覆盖率和百万级规模验证证明它能用、
可信、跑得动。

\vspace{1.1cm}
{\color{IRTCgold}\rule{2.6cm}{1.2pt}}\par
\vspace{0.3cm}
{\HeadFont\fontsize{12}{17}\selectfont\color{IRTCink}马崑翔}\par
{\LightFont\fontsize{10}{15}\selectfont\color{IRTCslate}%
 IRTC 创建人 / 维护人 \quad · \quad 2026 年 8 月}\par

\end{document}
